% ============================================================
%  PR\'EFACE — Syst\`emes Dynamiques et Chaos
% ============================================================
\chapter*{Pr\'eface}
\addcontentsline{toc}{chapter}{Pr\'eface}

\section*{Motivation et objectifs}

La th\'eorie des syst\`emes dynamiques constitue l'un des piliers fondamentaux
des math\'ematiques appliqu\'ees modernes. N\'ee des travaux pionniers de
Henri Poincar\'e sur le probl\`eme des trois corps \`a la fin du
\textsc{xix}\textsuperscript{e} si\`ecle, cette discipline a connu un
essor spectaculaire au cours du \textsc{xx}\textsuperscript{e} si\`ecle,
notamment gr\^ace aux contributions d'Alexandre Lyapounov, George Birkhoff,
Andre\"i Kolmogorov, Stephen Smale et Edward Lorenz.

Ce cours s'adresse aux \'etudiants de niveau Master (M1--M2) en
math\'ematiques, physique th\'eorique ou ing\'enierie, poss\'edant des bases
solides en analyse r\'eelle, alg\`ebre lin\'eaire et \'equations
diff\'erentielles ordinaires. Notre objectif est triple :

\begin{enumerate}
    \item \textbf{Ma\^itriser les fondements th\'eoriques} : existence et
    unicit\'e des solutions, stabilit\'e au sens de Lyapounov, th\'eorie
    spectrale des syst\`emes lin\'eaires, vari\'et\'es centrales et r\'eduction
    dimensionnelle.

    \item \textbf{Comprendre les ph\'enom\`enes qualitatifs} : cycles limites,
    bifurcations, chaos d\'eterministe, attracteurs \'etranges et
    comportement ergodique.

    \item \textbf{D\'evelopper des comp\'etences computationnelles} :
    simulation num\'erique en Python, calcul des exposants de Lyapounov,
    diagrammes de bifurcation et visualisation des attracteurs.
\end{enumerate}

\section*{Organisation du cours}

Le cours est structur\'e en onze chapitres, organis\'es selon une progression
logique allant des d\'efinitions fondamentales vers les concepts les plus
avanc\'es.

\medskip

\noindent\textbf{Partie I --- Fondements (Chapitres 1--3).}
Le chapitre~1 pose les d\'efinitions de base : syst\`emes dynamiques continus
et discrets, flots, orbites, ensembles limites. Le chapitre~2 introduit la
notion de stabilit\'e au sens de Lyapounov et les m\'ethodes directes et
indirectes. Le chapitre~3 traite la classification compl\`ete des portraits
de phase des syst\`emes lin\'eaires plans.

\medskip

\noindent\textbf{Partie II --- Syst\`emes non lin\'eaires (Chapitres 4--6).}
Le chapitre~4 d\'eveloppe les techniques de lin\'earisation et la th\'eorie
des vari\'et\'es centrales. Le chapitre~5 \'etudie les cycles limites et le
th\'eor\`eme de Poincar\'e-Bendixson. Le chapitre~6 pr\'esente la th\'eorie
des bifurcations : n\oe ud-col, fourche, Hopf.

\medskip

\noindent\textbf{Partie III --- Chaos et ergodicit\'e (Chapitres 7--11).}
Le chapitre~7 introduit les d\'efinitions formelles du chaos. Le chapitre~8
traite des exposants de Lyapounov. Le chapitre~9 \'etudie les attracteurs
\'etranges, notamment ceux de Lorenz et H\'enon. Le chapitre~10 aborde les
syst\`emes dynamiques discrets. Enfin, le chapitre~11 offre une introduction
\`a la th\'eorie ergodique.

\section*{Pr\'erequis}

Les pr\'erequis essentiels pour ce cours sont :

\begin{itemize}
    \item \textbf{Analyse r\'eelle} : topologie de $\R^n$, convergence
    uniforme, th\'eor\`emes de points fixes, int\'egrales d\'ependant d'un
    param\`etre.

    \item \textbf{Alg\`ebre lin\'eaire} : valeurs propres, d\'ecomposition
    de Jordan, exponentielle de matrice, formes bilin\'eaires.

    \item \textbf{\'Equations diff\'erentielles ordinaires} : th\'eor\`eme de
    Cauchy-Lipschitz, solutions maximales, lemme de Gr\"onwall.

    \item \textbf{Th\'eorie de la mesure} (pour le chapitre~11) : mesures
    de Borel, int\'egration de Lebesgue, th\'eor\`emes de convergence.

    \item \textbf{Programmation Python} (pour les travaux pratiques) :
    biblioth\`eques \texttt{numpy}, \texttt{scipy}, \texttt{matplotlib}.
\end{itemize}

\section*{Conventions et notations}

Tout au long de ce cours, nous adoptons les conventions suivantes :

\begin{itemize}
    \item $\R$, $\C$, $\N$, $\Z$ d\'esignent respectivement les ensembles
    des r\'eels, complexes, entiers naturels et entiers relatifs.

    \item $\norm{\cdot}$ d\'esigne une norme sur $\R^n$ (g\'en\'eralement la
    norme euclidienne $\norm{x} = \sqrt{x_1^2 + \cdots + x_n^2}$).

    \item $B(x_0, r) = \{x \in \R^n : \norm{x - x_0} < r\}$ est la boule
    ouverte de centre $x_0$ et de rayon $r$.

    \item $\dot{x}$ d\'esigne $\frac{dx}{dt}$, la d\'eriv\'ee par rapport au
    temps.

    \item $Df(x)$ ou $J_f(x)$ d\'esigne la matrice jacobienne de $f$ au
    point $x$.

    \item $\sigma(A)$ d\'esigne le spectre (ensemble des valeurs propres) de
    la matrice $A$.

    \item $\text{Re}(\lambda)$ et $\text{Im}(\lambda)$ d\'esignent les parties
    r\'eelle et imaginaire de $\lambda \in \C$.

    \item Les th\'eor\`emes, d\'efinitions et propositions sont num\'erot\'es
    par chapitre.
\end{itemize}

\section*{Perspective historique}

La th\'eorie des syst\`emes dynamiques poss\`ede une histoire riche et
fascinante que nous r\'esumons bri\`evement.

\paragraph{Les origines (1880--1920).}
Henri Poincar\'e, dans son \oe uvre monumentale \emph{Les m\'ethodes
nouvelles de la m\'ecanique c\'eleste} (1892--1899), a pos\'e les fondations
de l'approche qualitative des \'equations diff\'erentielles. Il a introduit
les notions de portrait de phase, de section de Poincar\'e et a entrevu la
possibilit\'e d'un comportement chaotique dans les syst\`emes d\'eterministes.

\paragraph{La stabilit\'e (1892--1950).}
Alexandre Lyapounov a d\'evelopp\'e sa th\'eorie de la stabilit\'e dans sa
th\`ese de 1892, \emph{Probl\`eme g\'en\'eral de la stabilit\'e du
mouvement}. Ses m\'ethodes, tant directes qu'indirectes, restent parmi les
outils les plus puissants de la th\'eorie moderne.

\paragraph{L'\'ecole topologique (1920--1960).}
George Birkhoff, puis les math\'ematiciens sovi\'etiques (Andronov,
Pontriaguine, Kolmogorov), ont d\'evelopp\'e l'approche topologique et
ergodique. Le th\'eor\`eme KAM (Kolmogorov-Arnold-Moser, 1954--1963)
a r\'esolu partiellement le probl\`eme de la stabilit\'e des syst\`emes
hamiltoniens.

\paragraph{La r\'evolution du chaos (1963--pr\'esent).}
Edward Lorenz a d\'ecouvert en 1963 la sensibilit\'e aux conditions initiales
dans un mod\`ele m\'et\'eorologique simplifi\'e. Les travaux de Smale
(fer \`a cheval), Ruelle et Takens (attracteurs \'etranges), Feigenbaum
(universalit\'e), et May (dynamique des populations) ont transform\'e notre
compr\'ehension des syst\`emes d\'eterministes.

\section*{Applications}

Les syst\`emes dynamiques trouvent des applications dans de nombreux domaines :

\begin{itemize}
    \item \textbf{Physique} : m\'ecanique c\'eleste, m\'ecanique des fluides
    (turbulence), physique des lasers, circuits \'electroniques non lin\'eaires.

    \item \textbf{Biologie} : dynamique des populations (mod\`eles
    proie-pr\'edateur), \'epid\'emiologie (mod\`ele SIR), neurosciences
    (mod\`ele de Hodgkin-Huxley), rythmes circadiens.

    \item \textbf{Chimie} : r\'eactions oscillantes (Belousov-Zhabotinsky),
    cin\'etique enzymatique.

    \item \textbf{\'Economie} : mod\`eles de croissance, cycles \'economiques,
    dynamique des march\'es financiers.

    \item \textbf{Ing\'enierie} : contr\^ole des syst\`emes, robotique,
    a\'eronautique (flutter), r\'eseaux \'electriques.

    \item \textbf{Climatologie} : mod\`eles climatiques, pr\'evision
    m\'et\'eorologique, El Ni\~no.
\end{itemize}

\section*{Comment utiliser ce cours}

Chaque chapitre contient :
\begin{itemize}
    \item Des \textbf{d\'efinitions} et \textbf{th\'eor\`emes} rigoureusement
    \'enonc\'es, avec d\'emonstrations d\'etaill\'ees.
    \item Des \textbf{exemples} illustrant les concepts abstraits.
    \item Des \textbf{remarques} apportant des compl\'ements ou des mises en
    garde.
    \item Des \textbf{exercices} de difficult\'e progressive.
    \item Des \textbf{simulations Python} permettant de visualiser les
    ph\'enom\`enes.
\end{itemize}

Nous encourageons vivement le lecteur \`a :
\begin{enumerate}
    \item Travailler les d\'emonstrations en d\'etail.
    \item R\'esoudre les exercices propos\'es.
    \item Ex\'ecuter et modifier les codes Python pour d\'evelopper une
    intuition num\'erique.
    \item \'Etablir des liens entre les diff\'erents chapitres.
\end{enumerate}

\section*{R\'ef\'erences principales}

\begin{enumerate}
    \item M.W.~Hirsch, S.~Smale, R.L.~Devaney, \emph{Differential Equations,
    Dynamical Systems, and an Introduction to Chaos}, 3\`eme \'ed.,
    Academic Press, 2012.
    \item S.H.~Strogatz, \emph{Nonlinear Dynamics and Chaos}, 2\`eme \'ed.,
    Westview Press, 2015.
    \item L.~Perko, \emph{Differential Equations and Dynamical Systems},
    3\`eme \'ed., Springer, 2001.
    \item A.~Katok, B.~Hasselblatt, \emph{Introduction to the Modern Theory
    of Dynamical Systems}, Cambridge University Press, 1995.
    \item J.~Guckenheimer, P.~Holmes, \emph{Nonlinear Oscillations, Dynamical
    Systems, and Bifurcations of Vector Fields}, Springer, 1983.
    \item K.T.~Alligood, T.D.~Sauer, J.A.~Yorke, \emph{Chaos: An Introduction
    to Dynamical Systems}, Springer, 1996.
\end{enumerate}

\section*{Remerciements}

Ce cours doit beaucoup aux ouvrages classiques cit\'es ci-dessus, ainsi
qu'aux nombreuses discussions avec des coll\`egues et \'etudiants. Nous
remercions \'egalement la communaut\'e math\'ematique francophone pour
son riche h\'eritage dans ce domaine, depuis Poincar\'e jusqu'\`a nos jours.

\vfill
\begin{flushright}
\textit{Bonne lecture et bon travail !}
\end{flushright}
